% Jeux de données et challenges
\section{Jeux de données et challenges}
\label{sec:datasets}

\subsection{Détection vue du drone : VisDrone et UAVDT}

\paragraph{VisDrone.} \citet{arxiv-2001.06303-detection-tracking-meet-drones-challenge}\pp{}
résument trois éditions du challenge (2018--2020). Le jeu compte 263 clips (179\,264
images) et 10\,209 images fixes, plus de 2,5 millions de boîtes sur 10 catégories, prises
par des DJI dans 14 villes chinoises, jusqu'à 2000$\times$1500 en image fixe
\src{§1}. La tâche DET se découpe en 6\,471 / 548 / 1\,580 / 1\,610 images
(train / val / test-challenge / \gl{testdev}{test-dev}) \src{§4.1} et suit le protocole \glslink{ap}{AP} \gl{COCO}{COCO}. Le
meilleur résultat de 2020, DroneEye2020 (Cascade R-CNN + DetectoRS), atteint 34,57 d'\glslink{ap}{AP}
\src{Table~2} ; le meilleur détecteur reste sous 25\,\% sur \emph{person} et
\emph{bicycle} \src{§4.5}. Les auteurs reconnaissent que les soumissions privilégient la
précision sur l'efficacité : les méthodes vidéo tournent à moins de 10\,\glslink{fps}{FPS} sur une \glslink{gtx1080ti}{GTX
1080Ti} \src{§5.3, §8}. Le §8 formule d'ailleurs le problème de la thèse : les ressources
d'un drone rendent impossible le déploiement de gros modèles.

\paragraph{UAVDT.} \citet{arxiv-1804.00518-unmanned-aerial-vehicle-benchmark-object}\pp{}
(version ECCV~2018 probable, non vérifiée) proposent 100 séquences urbaines à 1080$\times$540,
environ 80\,000 images annotées et 0,84\,M de boîtes de véhicules, avec des attributs
d'altitude, de vue et de météo \src{§2.1}. 27,5\,\% des objets font moins de 400 pixels
\src{§4}. Le meilleur détecteur (R-FCN) obtient 34,35\,\% d'\glslink{ap}{AP}, mais à un seuil d'\glslink{iou}{IoU} de
0,7 \src{§3.1} : ces chiffres ne se comparent pas aux \glslink{ap50}{AP@0.5} publiés ensuite sur UAVDT.
L'attribut altitude relie directement la précision à un paramètre de vol \deduit.

\subsection{Détection sous contrainte d'énergie : le DAC-SDC}

Le Low Power Object Detection Challenge du DAC est le seul contexte du corpus où la
précision, le débit et l'énergie sont mesurés par un tiers selon un protocole commun, sur
plateforme imposée. \citet{2021-xu-dac-sdc-low-power-object} décrivent l'édition 2018 :
images DJI 640$\times$360, un seul objet par image, la plupart occupant 1 à 2\,\% de
l'image ; plateformes \glslink{pynq-z1}{PYNQ-Z1} et \glslink{tx2}{Jetson TX2} ; énergie mesurée au wattmètre pour le \gl{FPGA}{FPGA}
\src{§2}. \citet{2024-zhang-comprehensive-analysis-dac-sdc-fpga} étendent l'analyse aux
éditions 2018--2022 (\glslink{pynq-z1}{PYNQ-Z1} puis \glslink{ultra96}{Ultra96}) et en tirent des règles de co-design
(section~\ref{sec:fpga}).

Deux réserves pèsent sur toute réutilisation. La métrique est un \glslink{iou-dacsdc}{IoU moyen} sans classe
sur un objet unique, qui ne se compare pas à une \glslink{map}{mAP} multi-classe
\citep[§2.3]{2024-zhang-comprehensive-analysis-dac-sdc-fpga}. Et la règle de mesure
change : à partir de 2022, le temps ne compte plus le chargement des images \src{§2.3},
ce qui explique une partie du saut de débit de cette édition.

\subsection{Contre-drone : des jeux surtout sol $\rightarrow$ air}

Le tableau~\ref{tab:datasets-cuav} classe les jeux du corpus par point de vue. Le constat
est net : les grands jeux (Anti-UAV, Anti-UAV410, Drone-vs-Bird, DUT Anti-UAV) sont pris
depuis le sol ou ne déclarent pas le capteur ; les jeux \gl{airair}{air $\rightarrow$ air} sont petits
et récents. Aucun ne comporte de piste « efficacité » (\glslink{fps}{FPS}, puissance, plateforme) : le
survey de \citet{arxiv-2507.10006-vision-anti-unmanned-aerial-technology}\pp{} le relève
lui-même \src{§V-A}.

\begin{table}[!htbp]
\centering\footnotesize
\caption{Jeux de contre-drone du corpus, par point de vue du capteur.}
\label{tab:datasets-cuav}
\begin{tabularx}{\textwidth}{@{}P{3.1cm}P{2.2cm}LP{2.4cm}@{}}
\toprule
Jeu, article & Point de vue & Contenu & Métrique \\
\midrule
Det-Fly \citep{2021-zheng-air-air-visual-detection-micro} & air $\rightarrow$ air &
  13\,271 images 4K, un DJI Mavic filmé par un M210, 10--100\,m ; public \src{§III} & \glslink{ap}{AP} \\
NPS-Drones, FL-Drones, ré-annotés \citep{arxiv-2103.17242-dogfight-detecting-drones-drones-videos}\pp &
  air $\rightarrow$ air & 50 et 14 vidéos ; cibles de 0,05 et 0,07\,\% de l'image \src{§1} & \glslink{ap}{AP}, \glslink{f1}{F1} \\
ARD-MAV \citep{arxiv-2312.11008-global-local-mav-detection-challenging}\pp & air $\rightarrow$ air &
  60 vidéos 1920$\times$1080, 106\,665 images, cibles de 0,02\,\% de l'image en moyenne
  \src{Table~I} & \glslink{ap50}{AP50} \\
Anti2 \citep{2024-gao-dataset-detection-method-unmanned-aerial} & mixte & 5\,062 images, 4 classes
  (drone, oiseau, hélicoptère, avion) \src{§3} & \glslink{map50}{mAP@0.5} \\
Anti-UAV \citep{arxiv-2101.08466-anti-uav-large-multi-modal}\pp & non déclaré (sol probable) &
  318 paires RGB--IR non alignées, cibles < 40\,px \src{Table~I} & \glslink{msa}{mSA} \\
Anti-UAV410 \citep{2024-huang-anti-uav410-thermal-infrared-benchmark} & non déclaré (sol probable) &
  410 séquences IR, plus de 100 à drone < 10\,px \src{§III} & \glslink{msa}{SA} \\
MAV-VID, Drone-vs-Bird \citep{arxiv-2103.13933-unmanned-aerial-vehicle-visual-detection}\pp &
  surtout \gl{solair}{sol $\rightarrow$ air} & 64 et 77 vidéos ; Drone-vs-Bird : cibles de 0,10\,\% \src{§3.1} & \glslink{ap50}{AP50} \\
Halmstad \citep{arxiv-2007.07396-real-time-drone-detection-tracking} & sol $\rightarrow$ air &
  650 vidéos IR + visible, 90 clips audio ; public \src{§IV} & \glslink{f1}{F1} \\
\bottomrule
\end{tabularx}
\end{table}

\paragraph{Suivi vu du drone.} WebUAV-3M
\citep{arxiv-2201.07425-webuav-3m-benchmark-unveiling-power}\pp{} compte 4\,500 vidéos et
3,3\,M d'images issues surtout de YouTube \src{§3.2}. Il évalue 43 trackers sur \gl{GPU}{GPU} de
bureau et montre que la précision s'effondre sous 30 pixels de cible \src{§5.6}.

\paragraph{Ce que les jeux ne mesurent pas.} Aucun des grands benchmarks du corpus
(VisDrone, UAVDT, Anti-UAV, Anti-UAV410, WebUAV-3M, Det-Fly) ne rapporte de débit sur
plateforme embarquée ni de puissance. Seul le \glslink{dac-sdc}{DAC-SDC} le fait, au prix d'une tâche
simplifiée. Les \glslink{ap}{AP} par taille d'objet sont rares même dans les jeux dédiés aux petits
objets : VisDrone ne publie pas d'\glslink{aps}{AP$_S$} séparé dans ses tableaux de synthèse
\citep[§4]{arxiv-2001.06303-detection-tracking-meet-drones-challenge}.
